\documentclass[a4paper,12pt]{article}
\usepackage[utf8]{inputenc}
\usepackage[T2A]{fontenc}
\usepackage[russian,english]{babel}
\usepackage{geometry}
\geometry{left=2.5cm,right=1.5cm,top=2cm,bottom=2cm}
\usepackage{graphicx}
\usepackage{hyperref}
\usepackage{verbatim}
\usepackage{indentfirst}
\begin{document}
% --- TITLE PAGE ---
\begin{titlepage}
\centering
{\large Министерство науки и высшего образования Российской Федерации\par}
\vspace{0.5cm}
{\large [Университет ИТМО]\par}
\vspace{2cm}
{\Large Отчёт по лабораторной работе\par}
\vspace{0.5cm}
{\large Знакомство с системой контроля версий Git и сервисом GitHub\par}
\vspace{2cm}
\begin{flushright}
Выполнил: Каира Еммануэль\\
Группа: R 3343\\
Проверил: Сениченкова Яна Олеговна
\end{flushright}
\vfill
{\large [Ваш город], 2026\par}
\end{titlepage}
\tableofcontents
\newpage
% --- TASK TEXT ---
\section{Текст задания}
\textbf{Цель:} поразмышлять о своих целях и задачах на текущий семестр. Создать файлы в репозитории, освоить работу с Git из командной строки, подготовить отчет.
\textbf{Задание 1:} Создать файл goals.md, добавить заголовки, текст о своем опыте, нумерованный и ненумерованный списки, создать папку resources с изображением.
\textbf{Задание 2:} Клонировать репозиторий, создать info.md, выполнить коммиты, создать .gitignore и объяснить результат игнорирования файла data/dataset.csv.
\textbf{Задание 3:} Подготовить отчет в формате PDF и загрузить его в репозиторий.
% --- TASK 1 ---
\section{Результат выполнения Задания 1}
\subsection{Содержимое файла goals.md}
\begin{verbatim}
% 
# Мои цели



1\. Научиться читать чужой код и работать в команде.

2.Собрать портфолио проектов для резюме. 

3\.Научиться читать чужой код и работать в команде.

## Задачи



\- Изучить основные команды Git: `add`, `commit`, `push`, `pull`.

\- Понять, как работает `.gitignore`.

\- Подготовить отчёт в LaTeX и загрузить его в репозиторий.

\- Научиться форматировать документы в Markdown.

\- Регулярно делать коммиты с понятными сообщениями.



!\[Иллюстрация](resources/image.png)


 goals.md
% (Скопируйте текст из Блокнота и вставьте сюда)
\end{verbatim}
\textit{Примечание: Изображение было успешно загружено в папку resources и добавлено в файл goals.md.}
% --- TASK 2 ---
\section{Результат выполнения Задания 2}
\subsection{Содержимое файла info.md}
\begin{verbatim}
% Каира Еммануель

windows 11 Home

17/09/2026 17:25

 info.md
% (3 строки: Каира Еммануель, windows 11 Home, 17/09/2026 17:25)
\end{verbatim}
\subsection{Содержимое файла .gitignore}
\begin{verbatim}
data/
\end{verbatim}
\subsection{Объяснение результата попытки закоммитить файл data/dataset.csv}
В файл .gitignore была добавлена строка \texttt{data/}. Это указывает Git игнорировать всё содержимое папки data. Поэтому при выполнении команды \texttt{git add data/dataset.csv} Git проигнорировал этот файл, и он не был добавлен в индекс. Соответственно, коммит с этим файлом не создался. Это доказывает, что .gitignore работает корректно.
% --- CONCLUSION ---
\section{Вывод}
В ходе выполнения данной лабораторной работы я ознакомился с системой контроля версий Git и сервисом GitHub. Я научился создавать репозитории, работать с файлами через веб-интерфейс и командную строку, а также освоил базовые команды Git (add, commit, push) и настройку .gitignore.
\end{document}